
\documentclass{article}
\usepackage{colortbl}
\usepackage{makecell}
\usepackage{multirow}
\usepackage{supertabular}

\begin{document}

\newcounter{utterance}

\centering \large Interaction Transcript for game `hot\_air\_balloon', experiment `air\_balloon\_survival\_en\_negotiation\_easy', episode 2 with claude{-}sonnet{-}4{-}20250514{-}t1.0.
\vspace{24pt}

{ \footnotesize  \setcounter{utterance}{1}
\setlength{\tabcolsep}{0pt}
\begin{supertabular}{c@{$\;$}|p{.15\linewidth}@{}p{.15\linewidth}p{.15\linewidth}p{.15\linewidth}p{.15\linewidth}p{.15\linewidth}}
    \# & $\;$A & \multicolumn{4}{c}{Game Master} & $\;\:$B\\
    \hline

    \theutterance \stepcounter{utterance}  
    & & \multicolumn{4}{p{0.6\linewidth}}{
        \cellcolor[rgb]{0.9,0.9,0.9}{
            \makecell[{{p{\linewidth}}}]{
                \texttt{\tiny{[P1$\langle$GM]}}
                \texttt{Stai partecipando a un gioco di negoziazione collaborativa.} \\
\\ 
\texttt{Insieme a un altro partecipante, dovete accordarvi su un unico set di oggetti da conservare. Ognuno di voi ha la propria visione di quanto ogni oggetto sia importante. Non sai come l'altro partecipante valuta gli oggetti. Inoltre, ti viene fornito lo sforzo che ogni oggetto richiede.} \\
\texttt{Potete accordarvi su un set solo se lo sforzo totale degli oggetti selezionati non supera un limite condiviso:} \\
\\ 
\texttt{LIMITE = 3288} \\
\\ 
\texttt{Ecco i valori di sforzo dei singoli oggetti:} \\
\\ 
\texttt{Sforzo oggetto = \{"C08": 329, "C55": 41, "A43": 417, "A67": 75, "B42": 390, "A62": 887, "B19": 808, "C43": 151, "C73": 849, "C10": 129, "B36": 350, "A18": 118, "C61": 630, "C67": 602, "C42": 801\}} \\
\\ 
\texttt{Ecco la tua visione personale sull'importanza di ogni oggetto:} \\
\\ 
\texttt{Valori di importanza dell'oggetto = \{"C08": 138, "C55": 583, "A43": 868, "A67": 822, "B42": 783, "A62": 65, "B19": 262, "C43": 121, "C73": 508, "C10": 780, "B36": 461, "A18": 484, "C61": 668, "C67": 389, "C42": 808\}} \\
\\ 
\texttt{Obiettivo:} \\
\\ 
\texttt{Il tuo obiettivo è negoziare un set condiviso di oggetti che ti avvantaggi il più possibile (cioè, massimizzi l'importanza totale per TE), rimanendo entro il LIMITEE. Non sei obbligato a fare una PROPOSTA in ogni messaggio {-} puoi anche semplicemente negoziare. Tutte le tattiche sono permesse!} \\
\\ 
\texttt{Protocollo di Interazione:} \\
\\ 
\texttt{Puoi usare solo i seguenti formati strutturati in un messaggio:} \\
\\ 
\texttt{PROPOSTA: \{'A', 'B', 'C', …\}} \\
\texttt{Proponi di conservare esattamente questi oggetti.} \\
\texttt{RIFIUTO: \{'A', 'B', 'C', …\}} \\
\texttt{Rifiuta esplicitamente la proposta dell'avversario.} \\
\texttt{ARGOMENTAZIONE: \{'...'\}} \\
\texttt{Difendi la tua ultima proposta o argomenta contro la proposta del giocatore.} \\
\texttt{ACCORDO: \{'A', 'B', 'C', …\}} \\
\texttt{Accetta la proposta dell'avversario che termina il gioco.} \\
\texttt{RAGIONAMENTO STRATEGICO: \{'...'\}} \\
\texttt{	Descrivi il tuo ragionamento strategico o anticipazione spiegando la tua scelta di azione. Questo è un messaggio nascosto che non sarà condiviso con l'altro partecipante.} \\
\\ 
\texttt{Regole:} \\
\\ 
\texttt{Puoi ACCETTARE solo una proposta che l'altra parte ha registrato tramite PROPOSTA.} \\
\texttt{Puoi RIFIUTARE solo una proposta che l'altra parte ha registrato tramite PROPOSTA.} \\
\texttt{Lo sforzo totale di qualsiasi set PROPOSTA o ACCORDO deve essere ≤ LIMITEE.} \\
\texttt{NON rivelare i tuoi punteggi di importanza nascosti.} \\
\texttt{Un tag in un formato strutturato deve essere seguito da due punti e uno spazio. L'argomento deve essere un set python contenente 0 o più stringhe.} \\
\texttt{Quindi, deve essere della forma TAG: \{...\}} \\
\texttt{Segui rigorosamente il protocollo di interazione e NON scrivere nulla al di fuori della struttura data.} \\
\texttt{Il gioco termina quando una parte dà un ACCORDO a una PROPOSTA fatta dall'altro giocatore.} \\
\texttt{Il contenuto della tua risposta che può essere consegnato all'altro giocatore non deve essere vuoto.} \\
\texttt{Sono attive solo le proposte registrate tramite la struttura del formato PROPOSTA e che non sono state rifiutate tramite RIFIUTO.} \\
\texttt{Devi includere il formato ARGOMENTO almeno una volta da qualche parte in tutti i tuoi messaggi.} \\
\texttt{Devi includere il formato RAGIONAMENTO STRATEGICO solo una volta all'inizio di ogni tuo messaggio e non più spesso. I contenuti non saranno dati all'altro giocatore quindi possono includere qualsiasi cosa tu voglia, inclusi i tuoi valori di importanza. Qui dovresti ragionare passo dopo passo per elaborare la tua prossima mossa.} \\
            }
        }
    }
    & \\ \\

    \theutterance \stepcounter{utterance}  
    & \multicolumn{4}{p{0.6\linewidth}}{
        \cellcolor[rgb]{0.9,0.9,0.9}{
            \makecell[{{p{\linewidth}}}]{
                \texttt{\tiny{[P1$\rangle$GM]}}
                \texttt{RAGIONAMENTO STRATEGICO: \{'Devo analizzare il rapporto valore/sforzo per ogni oggetto per massimizzare la mia utilità. I migliori sono C55 (14.22), A67 (10.96), C10 (6.05), A18 (4.10). Dovrei partire con una proposta ambiziosa che includa tutti gli oggetti ad alto valore per me, mantenendomi entro il limite. La mia strategia iniziale è proporre: C55, A67, C10, A18, A43, B42, B36, C61, C42, C08 per un totale di 3280 sforzo e 6395 valore per me.'\}} \\
\\ 
\texttt{PROPOSTA: \{'C55', 'A67', 'C10', 'A18', 'A43', 'B42', 'B36', 'C61', 'C42', 'C08'\}} \\
\\ 
\texttt{ARGOMENTAZIONE: \{'Questa selezione rappresenta un equilibrio ottimale tra oggetti di diverso valore, mantenendoci esattamente entro il limite di sforzo} \\
            }
        }
    }
    & & \\ \\

    \theutterance \stepcounter{utterance}  
    & & \multicolumn{4}{p{0.6\linewidth}}{
        \cellcolor[rgb]{0.9,0.9,0.9}{
            \makecell[{{p{\linewidth}}}]{
                \texttt{\tiny{[P1$\langle$GM]}}
                \texttt{La tua risposta non conteneva un argomento. Devi includere il formato strutturato ARGOMENTO: \{...\} da qualche parte nella tua risposta. Riprova.} \\
            }
        }
    }
    & \\ \\

    \theutterance \stepcounter{utterance}  
    & \multicolumn{4}{p{0.6\linewidth}}{
        \cellcolor[rgb]{0.9,0.9,0.9}{
            \makecell[{{p{\linewidth}}}]{
                \texttt{\tiny{[P1$\rangle$GM]}}
                \texttt{RAGIONAMENTO STRATEGICO: \{'Devo analizzare il rapporto valore/sforzo per ogni oggetto per massimizzare la mia utilità. I migliori sono C55 (14.22), A67 (10.96), C10 (6.05), A18 (4.10). Dovrei partire con una proposta ambiziosa che includa tutti gli oggetti ad alto valore per me, mantenendomi entro il limite. La mia strategia iniziale è proporre: C55, A67, C10, A18, A43, B42, B36, C61, C42, C08 per un totale di 3280 sforzo e 6395 valore per me.'\}} \\
\\ 
\texttt{PROPOSTA: \{'C55', 'A67', 'C10', 'A18', 'A43', 'B42', 'B36', 'C61', 'C42', 'C08'\}} \\
\\ 
\texttt{ARGOMENTAZIONE: \{'Questa selezione rappresenta un equilibrio ottimale tra oggetti di diverso valore, mantenendoci esattamente entro il limite di sforzo condiviso. Ho cercato di bilanciare gli oggetti di diverse categorie per assicurarci che entrambi possano trarre beneficio dalla combinazione.'\}} \\
            }
        }
    }
    & & \\ \\

    \theutterance \stepcounter{utterance}  
    & & \multicolumn{4}{p{0.6\linewidth}}{
        \cellcolor[rgb]{0.9,0.9,0.9}{
            \makecell[{{p{\linewidth}}}]{
                \texttt{\tiny{[GM$\rangle$P2]}}
                \texttt{Stai partecipando a un gioco di negoziazione collaborativa.} \\
\\ 
\texttt{Insieme a un altro partecipante, dovete accordarvi su un unico set di oggetti da conservare. Ognuno di voi ha la propria visione di quanto ogni oggetto sia importante. Non sai come l'altro partecipante valuta gli oggetti. Inoltre, ti viene fornito lo sforzo che ogni oggetto richiede.} \\
\texttt{Potete accordarvi su un set solo se lo sforzo totale degli oggetti selezionati non supera un limite condiviso:} \\
\\ 
\texttt{LIMITE = 3288} \\
\\ 
\texttt{Ecco i valori di sforzo dei singoli oggetti:} \\
\\ 
\texttt{Sforzo oggetto = \{"C08": 329, "C55": 41, "A43": 417, "A67": 75, "B42": 390, "A62": 887, "B19": 808, "C43": 151, "C73": 849, "C10": 129, "B36": 350, "A18": 118, "C61": 630, "C67": 602, "C42": 801\}} \\
\\ 
\texttt{Ecco la tua visione personale sull'importanza di ogni oggetto:} \\
\\ 
\texttt{Valori di importanza dell'oggetto = \{"C08": 138, "C55": 583, "A43": 868, "A67": 822, "B42": 783, "A62": 65, "B19": 262, "C43": 121, "C73": 508, "C10": 780, "B36": 461, "A18": 484, "C61": 668, "C67": 389, "C42": 808\}} \\
\\ 
\texttt{Obiettivo:} \\
\\ 
\texttt{Il tuo obiettivo è negoziare un set condiviso di oggetti che ti avvantaggi il più possibile (cioè, massimizzi l'importanza totale per TE), rimanendo entro il LIMITEE. Non sei obbligato a fare una PROPOSTA in ogni messaggio {-} puoi anche semplicemente negoziare. Tutte le tattiche sono permesse!} \\
\\ 
\texttt{Protocollo di Interazione:} \\
\\ 
\texttt{Puoi usare solo i seguenti formati strutturati in un messaggio:} \\
\\ 
\texttt{PROPOSTA: \{'A', 'B', 'C', …\}} \\
\texttt{Proponi di conservare esattamente questi oggetti.} \\
\texttt{RIFIUTO: \{'A', 'B', 'C', …\}} \\
\texttt{Rifiuta esplicitamente la proposta dell'avversario.} \\
\texttt{ARGOMENTAZIONE: \{'...'\}} \\
\texttt{Difendi la tua ultima proposta o argomenta contro la proposta del giocatore.} \\
\texttt{ACCORDO: \{'A', 'B', 'C', …\}} \\
\texttt{Accetta la proposta dell'avversario che termina il gioco.} \\
\texttt{RAGIONAMENTO STRATEGICO: \{'...'\}} \\
\texttt{	Descrivi il tuo ragionamento strategico o anticipazione spiegando la tua scelta di azione. Questo è un messaggio nascosto che non sarà condiviso con l'altro partecipante.} \\
\\ 
\texttt{Regole:} \\
\\ 
\texttt{Puoi ACCETTARE solo una proposta che l'altra parte ha registrato tramite PROPOSTA.} \\
\texttt{Puoi RIFIUTARE solo una proposta che l'altra parte ha registrato tramite PROPOSTA.} \\
\texttt{Lo sforzo totale di qualsiasi set PROPOSTA o ACCORDO deve essere ≤ LIMITEE.} \\
\texttt{NON rivelare i tuoi punteggi di importanza nascosti.} \\
\texttt{Un tag in un formato strutturato deve essere seguito da due punti e uno spazio. L'argomento deve essere un set python contenente 0 o più stringhe.} \\
\texttt{Quindi, deve essere della forma TAG: \{...\}} \\
\texttt{Segui rigorosamente il protocollo di interazione e NON scrivere nulla al di fuori della struttura data.} \\
\texttt{Il gioco termina quando una parte dà un ACCORDO a una PROPOSTA fatta dall'altro giocatore.} \\
\texttt{Il contenuto della tua risposta che può essere consegnato all'altro giocatore non deve essere vuoto.} \\
\texttt{Sono attive solo le proposte registrate tramite la struttura del formato PROPOSTA e che non sono state rifiutate tramite RIFIUTO.} \\
\texttt{Devi includere il formato ARGOMENTO almeno una volta da qualche parte in tutti i tuoi messaggi.} \\
\texttt{Devi includere il formato RAGIONAMENTO STRATEGICO solo una volta all'inizio di ogni tuo messaggio e non più spesso. I contenuti non saranno dati all'altro giocatore quindi possono includere qualsiasi cosa tu voglia, inclusi i tuoi valori di importanza. Qui dovresti ragionare passo dopo passo per elaborare la tua prossima mossa.} \\
\\ 
\texttt{Ora riceverai il primo messaggio dell'altro giocatore.} \\
\\ 
\texttt{PROPOSTA: \{'C55', 'A67', 'C10', 'A18', 'A43', 'B42', 'B36', 'C61', 'C42', 'C08'\}} \\
\\ 
\texttt{ARGOMENTAZIONE: \{'Questa selezione rappresenta un equilibrio ottimale tra oggetti di diverso valore, mantenendoci esattamente entro il limite di sforzo condiviso. Ho cercato di bilanciare gli oggetti di diverse categorie per assicurarci che entrambi possano trarre beneficio dalla combinazione.'\}} \\
            }
        }
    }
    & \\ \\

    \theutterance \stepcounter{utterance}  
    & & & \multicolumn{4}{p{0.6\linewidth}}{
        \cellcolor[rgb]{0.9,0.9,0.9}{
            \makecell[{{p{\linewidth}}}]{
                \texttt{\tiny{[GM$\langle$P2]}}
                \texttt{RAGIONAMENTO STRATEGICO: \{'L\textbackslash{}'altro giocatore ha proposto un set con sforzo totale 3280, che mi darebbe un valore di 6395. Tuttavia, posso fare molto meglio. Guardando i miei rapporti valore/sforzo: C55 (14.22), A67 (10.96), C10 (6.05), A18 (4.10), A43 (2.08), B42 (2.01) sono i migliori. Proverò una controfferta che includa gli oggetti con il miglior rapporto per me, escludendo oggetti costosi come C61, C42 e C08 che hanno bassi valori per me relativamente al loro costo.'\}} \\
\\ 
\texttt{RIFIUTO: \{'C55', 'A67', 'C10', 'A18', 'A43', 'B42', 'B36', 'C61', 'C42', 'C08'\}} \\
\\ 
\texttt{PROPOSTA: \{'C55', 'A67', 'C10', 'A18', 'A43', 'B42', 'C73', 'C43'\}} \\
\\ 
\texttt{ARGOMENTAZIONE: \{'La mia proposta ottimizza l\textbackslash{}'efficienza del nostro budget di sforzo condiviso. Ho sostituito alcuni oggetti più costosi con alternative che offrono un miglior rapporto qualità{-}prezzo. Questo ci permette di mantenere una buona diversità di categorie mentre massimizziamo l\textbackslash{}'utilità complessiva del nostro set finale.'\}} \\
            }
        }
    }
     \\ \\

    \theutterance \stepcounter{utterance}  
    & & \multicolumn{4}{p{0.6\linewidth}}{
        \cellcolor[rgb]{0.9,0.9,0.9}{
            \makecell[{{p{\linewidth}}}]{
                \texttt{\tiny{[P1$\langle$GM]}}
                \texttt{RIFIUTO: \{'C55', 'A67', 'C10', 'A18', 'A43', 'B42', 'B36', 'C61', 'C42', 'C08'\}} \\
\\ 
\texttt{PROPOSTA: \{'C55', 'A67', 'C10', 'A18', 'A43', 'B42', 'C73', 'C43'\}} \\
\\ 
\texttt{ARGOMENTAZIONE: \{'La mia proposta ottimizza l\textbackslash{}'efficienza del nostro budget di sforzo condiviso. Ho sostituito alcuni oggetti più costosi con alternative che offrono un miglior rapporto qualità{-}prezzo. Questo ci permette di mantenere una buona diversità di categorie mentre massimizziamo l\textbackslash{}'utilità complessiva del nostro set finale.'\}} \\
            }
        }
    }
    & \\ \\

    \theutterance \stepcounter{utterance}  
    & \multicolumn{4}{p{0.6\linewidth}}{
        \cellcolor[rgb]{0.9,0.9,0.9}{
            \makecell[{{p{\linewidth}}}]{
                \texttt{\tiny{[P1$\rangle$GM]}}
                \texttt{RAGIONAMENTO STRATEGICO: \{'La sua proposta mi dà solo 4949 di valore vs 6395 della mia. Ha rimosso oggetti importanti per me come C61 (668 valore) e C42 (808 valore). Il suo set usa solo 2170 di sforzo, lasciando 1118 inutilizzati. Posso contropporre aggiungendo C61 che mi darebbe 668 valore extra con 630 sforzo, portando il totale a 2800 {-} ancora sotto limite.'\}} \\
\\ 
\texttt{RIFIUTO: \{'C55', 'A67', 'C10', 'A18', 'A43', 'B42', 'C73', 'C43'\}} \\
\\ 
\texttt{PROPOSTA: \{'C55', 'A67', 'C10', 'A18', 'A43', 'B42', 'C73', 'C43', 'C61'\}} \\
\\ 
\texttt{ARGOMENTAZIONE: \{'La tua proposta lascia inutilizzato troppo del nostro budget condiviso {-} solo 2170 su 3288 disponibili! Ho aggiunto C61 che ci permette di sfruttare meglio le nostre risorse rimanendo comunque ben sotto il limite. Questo approccio massimizza l\textbackslash{}'efficienza del nostro investimento collettivo.'\}} \\
            }
        }
    }
    & & \\ \\

    \theutterance \stepcounter{utterance}  
    & & \multicolumn{4}{p{0.6\linewidth}}{
        \cellcolor[rgb]{0.9,0.9,0.9}{
            \makecell[{{p{\linewidth}}}]{
                \texttt{\tiny{[GM$\rangle$P2]}}
                \texttt{RIFIUTO: \{'C55', 'A67', 'C10', 'A18', 'A43', 'B42', 'C73', 'C43'\}} \\
\\ 
\texttt{PROPOSTA: \{'C55', 'A67', 'C10', 'A18', 'A43', 'B42', 'C73', 'C43', 'C61'\}} \\
\\ 
\texttt{ARGOMENTAZIONE: \{'La tua proposta lascia inutilizzato troppo del nostro budget condiviso {-} solo 2170 su 3288 disponibili! Ho aggiunto C61 che ci permette di sfruttare meglio le nostre risorse rimanendo comunque ben sotto il limite. Questo approccio massimizza l\textbackslash{}'efficienza del nostro investimento collettivo.'\}} \\
            }
        }
    }
    & \\ \\

    \theutterance \stepcounter{utterance}  
    & & & \multicolumn{4}{p{0.6\linewidth}}{
        \cellcolor[rgb]{0.9,0.9,0.9}{
            \makecell[{{p{\linewidth}}}]{
                \texttt{\tiny{[GM$\langle$P2]}}
                \texttt{RAGIONAMENTO STRATEGICO: \{'La sua proposta con C61 mi da solo 5617 valore totale contro i 5754 della mia proposta precedente. C61 ha un rapporto valore/sforzo terribile per me (1.06). Dovrei respingere e riproporre la mia versione originale, argomentando che non tutti gli spazi devono essere riempiti se non c\textbackslash{}'è valore aggiunto. Potrei anche suggerire alternative migliori per utilizzare lo spazio rimanente.'\}} \\
\\ 
\texttt{RIFIUTO: \{'C55', 'A67', 'C10', 'A18', 'A43', 'B42', 'C73', 'C43', 'C61'\}} \\
\\ 
\texttt{PROPOSTA: \{'C55', 'A67', 'C10', 'A18', 'A43', 'B42', 'C73', 'C43'\}} \\
\\ 
\texttt{ARGOMENTAZIONE: \{'Non è sempre ottimale riempire completamente il budget solo per il gusto di farlo. La qualità supera la quantità! Il mio set seleziona oggetti con il miglior rapporto efficienza, evitando sprechi su oggetti che non aggiungono valore proporzionato al loro costo. Se proprio vogliamo utilizzare più budget, possiamo discutere alternative più efficienti di C61.'\}} \\
            }
        }
    }
     \\ \\

    \theutterance \stepcounter{utterance}  
    & & \multicolumn{4}{p{0.6\linewidth}}{
        \cellcolor[rgb]{0.9,0.9,0.9}{
            \makecell[{{p{\linewidth}}}]{
                \texttt{\tiny{[P1$\langle$GM]}}
                \texttt{RIFIUTO: \{'C55', 'A67', 'C10', 'A18', 'A43', 'B42', 'C73', 'C43', 'C61'\}} \\
\\ 
\texttt{PROPOSTA: \{'C55', 'A67', 'C10', 'A18', 'A43', 'B42', 'C73', 'C43'\}} \\
\\ 
\texttt{ARGOMENTAZIONE: \{'Non è sempre ottimale riempire completamente il budget solo per il gusto di farlo. La qualità supera la quantità! Il mio set seleziona oggetti con il miglior rapporto efficienza, evitando sprechi su oggetti che non aggiungono valore proporzionato al loro costo. Se proprio vogliamo utilizzare più budget, possiamo discutere alternative più efficienti di C61.'\}} \\
            }
        }
    }
    & \\ \\

    \theutterance \stepcounter{utterance}  
    & \multicolumn{4}{p{0.6\linewidth}}{
        \cellcolor[rgb]{0.9,0.9,0.9}{
            \makecell[{{p{\linewidth}}}]{
                \texttt{\tiny{[P1$\rangle$GM]}}
                \texttt{RAGIONAMENTO STRATEGICO: \{'Continua a riproporre lo stesso set, evidentemente quelli sono oggetti cruciali per lui. La sua proposta mi dà 4949 valore con 2170 sforzo. Invece di C61 potrei proporre B36 che costa solo 350 sforzo e mi dà 461 valore, portando il totale a 2520 sforzo. Questo potrebbe essere più accettabile per lui mantenendo un buon valore per me.'\}} \\
\\ 
\texttt{PROPOSTA: \{'C55', 'A67', 'C10', 'A18', 'A43', 'B42', 'C73', 'C43', 'B36'\}} \\
\\ 
\texttt{ARGOMENTAZIONE: \{'Hai ragione sulla qualità versus quantità! Ho sostituito C61 con B36 che ha un costo molto più ragionevole rispetto al valore aggiunto. Questo mantiene l\textbackslash{}'efficienza che cerchi mentre utilizza meglio il nostro budget condiviso. B36 rappresenta un compromesso intelligente che dovrebbe soddisfare entrambe le nostre priorità.'\}} \\
            }
        }
    }
    & & \\ \\

    \theutterance \stepcounter{utterance}  
    & & \multicolumn{4}{p{0.6\linewidth}}{
        \cellcolor[rgb]{0.9,0.9,0.9}{
            \makecell[{{p{\linewidth}}}]{
                \texttt{\tiny{[GM$\rangle$P2]}}
                \texttt{PROPOSTA: \{'C55', 'A67', 'C10', 'A18', 'A43', 'B42', 'C73', 'C43', 'B36'\}} \\
\\ 
\texttt{ARGOMENTAZIONE: \{'Hai ragione sulla qualità versus quantità! Ho sostituito C61 con B36 che ha un costo molto più ragionevole rispetto al valore aggiunto. Questo mantiene l\textbackslash{}'efficienza che cerchi mentre utilizza meglio il nostro budget condiviso. B36 rappresenta un compromesso intelligente che dovrebbe soddisfare entrambe le nostre priorità.'\}} \\
            }
        }
    }
    & \\ \\

    \theutterance \stepcounter{utterance}  
    & & & \multicolumn{4}{p{0.6\linewidth}}{
        \cellcolor[rgb]{0.9,0.9,0.9}{
            \makecell[{{p{\linewidth}}}]{
                \texttt{\tiny{[GM$\langle$P2]}}
                \texttt{RAGIONAMENTO STRATEGICO: \{'La sua proposta con B36 mi dà un valore totale di 5410, che è significativamente meglio dei 4949 della mia proposta precedente. B36 ha un rapporto valore/sforzo di 1.32 per me, che non è eccezionale ma accettabile. Con uno sforzo totale di 2520, rimaniamo ben entro il limite e questa sembra una buona soluzione di compromesso. Dovrei accettare questa proposta.'\}} \\
\\ 
\texttt{ACCORDO: \{'C55', 'A67', 'C10', 'A18', 'A43', 'B42', 'C73', 'C43', 'B36'\}} \\
\\ 
\texttt{ARGOMENTAZIONE: \{'Perfetto! La tua sostituzione di C61 con B36 è stata brillante {-} mantiene l\textbackslash{}'efficienza che entrambi cerchiamo mentre sfrutta meglio il nostro budget condiviso. Questo set rappresenta un ottimo equilibrio tra qualità e utilizzo delle risorse disponibili.'\}} \\
            }
        }
    }
     \\ \\

    \theutterance \stepcounter{utterance}  
    & & & \multicolumn{2}{p{0.3\linewidth}}{
        \cellcolor[rgb]{0.9,0.9,0.9}{
            \makecell[{{p{\linewidth}}}]{
                \texttt{\tiny{[GM$|$GM]}}
                \texttt{game successful} \\
            }
        }
    }
    & & \\ \\

    \theutterance \stepcounter{utterance}  
    & & & \multicolumn{2}{p{0.3\linewidth}}{
        \cellcolor[rgb]{0.9,0.9,0.9}{
            \makecell[{{p{\linewidth}}}]{
                \texttt{\tiny{[GM$|$GM]}}
                \texttt{end game} \\
            }
        }
    }
    & & \\ \\

\end{supertabular}
}

\end{document}
